% Included by uat-professor-proof.tex. Not a standalone document.
\clearpage
\pdfbookmark[0]{پیوست‌های آموزشی}{bm-appendices}
\phantomsection\label{app:approx}
\pdfbookmark[1]{پیوست الف: زبان تقریب و دو قاعدهٔ اصلی}{bm-app-approx}
\ribbon{پیوست‌ها | اثبات اصلی در صفحهٔ \pageref{main:end} پایان یافته است}
\heading{پیوست الف. زبان تقریب و دو قاعدهٔ اصلی}
\textbf{محل استفاده:} گام ۱ و همهٔ حدگیری‌های اثبات.

\subhead{راهنمای پیوست‌ها}
ابزارهای فشردگی در پیوست ب؛ حساب دیفرانسیل و انتگرال در پ؛ درون‌یابی لاگرانژ در ت؛ هموارسازی در ث؛ مشتق نسبت به وزن در ج؛ وایرشتراس در چ؛ و استون–وایرشتراس در ح آمده‌اند. عنوان‌های PDF و ارجاع‌های داخل متن قابل کلیک‌اند.

\subhead{تقریب یکنواخت یعنی چه؟}
برای یک تابع پیوسته روی $K$، می‌نویسیم
\[
\normK{u}=\sup_{x\in K}|u(x)|.
\]
رابطهٔ $u_n\to u$ به‌طور یکنواخت یعنی این بزرگ‌ترین خطا به صفر برسد. در نتیجه برای هر دقت، یک شمارهٔ $n$ داریم که برای \emph{تمام} نقاط دامنه کافی است. در همگرایی نقطه‌ای، این شماره می‌تواند با نقطه عوض شود.

\subhead{تعریف دقیق مجموعهٔ قابل تقریب}
تابع پیوستهٔ $u$ را عضو $\V$ می‌نامیم اگر برای هر $\eta>0$، یک شبکهٔ متناهی با فعال‌سازی اصلی وجود داشته باشد که $\normK{u-F}<\eta$. در زبان آنالیز، این مجموعه «بستار یکنواختِ خانوادهٔ شبکه‌ها» است؛ برای استفاده از آن، دو قاعدهٔ زیر کافی‌اند.

\subhead{قاعدهٔ اول: ترکیب خطی}
اگر $u,v\in\V$ و $a,b\in\R$ باشند، $au+bv\in\V$. برای اثبات، $u$ و $v$ را با شبکه‌های $F,G$، هرکدام با خطای کمتر از $\eta/(1+|a|+|b|)$، تقریب می‌زنیم. آنگاه
\[
\normK{au+bv-(aF+bG)}
\leq |a|\normK{u-F}+|b|\normK{v-G}<\eta.
\]
تابع $aF+bG$ یک شبکه است؛ نورون‌های دو شبکه را کنار هم می‌گذاریم. با تکرار همین کار، حکم برای هر جمع متناهی برقرار می‌شود.

\subhead{قاعدهٔ دوم: حد یکنواخت}
اگر $u_n\in\V$ و $u_n\to u$ یکنواخت باشد، $u\in\V$. برای دقت $\eta$، ابتدا $n$ را طوری می‌گیریم که $\normK{u-u_n}<\eta/2$. سپس شبکه‌ای $F$ را با $\normK{u_n-F}<\eta/2$ انتخاب می‌کنیم. از نامساوی مثلثی، $\normK{u-F}<\eta$ نتیجه می‌شود. پس یک شبکهٔ \textbf{متناهی} کافی است؛ حدگیری، معماری نامتناهی تولید نکرده است.

\subhead{چرا تابع حد هنوز پیوسته است؟}
حد یکنواخت توابع پیوسته، پیوسته است. برای نقطهٔ $x_0$ و دقت $\eta$، یک $n$ با خطای یکنواخت کمتر از $\eta/3$ ثابت کنید. از پیوستگی $u_n$، برای $x$ نزدیک $x_0$، اختلاف $|u_n(x)-u_n(x_0)|<\eta/3$ است. با جمع سه خطا، $|u(x)-u(x_0)|<\eta$ به دست می‌آید.

\begin{intuition}
هیچ‌کدام از این دو قاعده نمی‌گوید حاصل‌ضرب دو تابعِ قابل تقریب نیز قابل تقریب است. در اثبات اصلی به چنین فرضی تکیه نکردیم.
\end{intuition}
\backtomain{step:one}

\appstart{ب}{فشردگی، کرانداری و پیوستگی یکنواخت}{app:compact}
\textbf{محل استفاده:} کنترل خطا برای همهٔ ورودی‌ها با یک انتخاب مشترک.

\subhead{قضیهٔ هاینه–بورل: صورت موردنیاز}
در $\R^d$، یک مجموعه فشرده است اگر و تنها اگر بسته و کراندار باشد. این صورت کلاسیکِ قضیه را به‌عنوان یک نتیجهٔ مقدماتی به کار می‌بریم. تعریف عمومی فشردگی این است: از هر پوششِ مجموعه با مجموعه‌های باز، می‌توان یک زیرپوشش متناهی انتخاب کرد.

برای مثال $[0,1]^d$ فشرده است. فشرده لازم نیست داخلِ ناتهی داشته باشد؛ یک دایرهٔ محیطی یا یک پاره‌خط نیز می‌تواند فشرده باشد. حاصل‌ضربِ تعداد متناهی فشردهٔ اقلیدسی هم فشرده است، چون بسته و کراندار می‌ماند.

\subhead{قضیهٔ بیشینه و کمینه}
اگر $u$ روی فشردهٔ ناتهی $K$ پیوسته باشد، مقدارهایش کراندارند و بیشینه و کمینه را می‌گیرد.

\textbf{دلیل.} تصویر پیوستهٔ فشرده، فشرده است: پوشش بازِ تصویر را با پیش‌تصویر به دامنه برمی‌گردانیم و زیرپوشش متناهی می‌گیریم. پس $u(K)$ در خط حقیقی بسته و کراندار است. سوپریمم و اینفیمم آن متناهی‌اند و، به دلیل بسته‌بودن، خودشان در مجموعه قرار دارند.

در اثبات اصلی، نتیجه می‌گیریم هر مختصهٔ $x_i$ و هر حاصل‌ضربِ متناهی از مختصات روی $K$ کراندار است. بنابراین این عامل‌ها نمی‌توانند یک خطای یکنواختِ رو به صفر را بی‌حد بزرگ کنند.

\subhead{قضیهٔ هاینه–کانتور}
هر تابع پیوسته روی مجموعهٔ فشردهٔ $K\subset\R^d$، یکنواختاً پیوسته است. یعنی برای هر $\eta>0$، عددی $\delta>0$ وجود دارد که برای \emph{تمام} $x,y\in K$،
\[
\|x-y\|<\delta\quad\Longrightarrow\quad |u(x)-u(y)|<\eta.
\]
\textbf{اثبات.} برای هر $z\in K$ شعاع $r_z>0$ انتخاب کنید که در فاصلهٔ کمتر از $2r_z$ از $z$، اختلاف مقدار تابع با $u(z)$ کمتر از $\eta/2$ باشد. توپ‌های $B(z,r_z)$ پوشش $K$ هستند. زیرپوشش متناهی با مراکز $z_1,\ldots,z_m$ انتخاب و $\delta=\min_i r_{z_i}$ تعریف کنید.

اگر $\|x-y\|<\delta$، نقطهٔ $x$ در یکی از توپ‌های انتخابی است. آنگاه فاصلهٔ $y$ از همان مرکز کمتر از $2r_{z_i}$ می‌شود. هر دو مقدار $u(x),u(y)$ کمتر از $\eta/2$ از $u(z_i)$ فاصله دارند؛ پس اختلافشان کمتر از $\eta$ است.

\begin{intuition}
پیوستگی معمولی برای هر نقطه یک اندازهٔ همسایگی می‌دهد؛ پیوستگی یکنواخت یک اندازهٔ مشترک برای همهٔ نقاط می‌دهد. این تفاوت، علتِ اصلی استفاده از دامنهٔ فشرده در مجموع‌های ریمان و تفاضل‌های مشتق است.
\end{intuition}
\backtomain{step:two}

\appstart{پ}{سه ابزار حساب دیفرانسیل و انتگرال}{app:calculus}
\textbf{محل استفاده:} هموارسازی، صفرشدن مشتق و ظاهرشدن توان‌های ورودی.

\subhead{قضیهٔ اساسی حساب ــ مشتقِ انتگرال}
اگر $v$ پیوسته و $I(t)=\int_a^t v(s)\,ds$ باشد، آنگاه $I'(t)=v(t)$ در نقاط داخلی بازه.

\textbf{دلیل.} برای $r\neq0$ داریم
\[
\frac{I(t+r)-I(t)}{r}-v(t)
=\frac1r\int_t^{t+r}\bigl(v(s)-v(t)\bigr)\,ds.
\]
قدر مطلق سمت راست از بیشترینِ $|v(s)-v(t)|$ بین $t$ و $t+r$ بیشتر نیست. این مقدار با پیوستگی به صفر می‌رود؛ برای $r<0$ نیز همین استدلال برقرار است.

\subhead{صورت دوم ــ انتگرالِ مشتق}
اگر $V'=v$ و $v$ پیوسته باشد، $\int_a^b v(s)\,ds=V(b)-V(a)$. زیرا مشتقِ $V-I$ صفر است و طبق نتیجهٔ مقدار میانگین، این اختلاف ثابت می‌ماند. نتیجهٔ ویژهٔ موردنیاز ما برای $g\in C^1$ چنین است:
\begin{equation}\label{eq:ftc-difference}
 g(z+a)-g(z)=a\int_0^1g'(z+\theta a)\,d\theta.
\end{equation}
این رابطه برای $a=0$ هم درست است. از آن برای تبدیل تفاضل‌ها به یک انتگرال با خطای قابل کنترل استفاده می‌کنیم.

\subhead{قاعدهٔ زنجیره‌ای}
اگر $u$ و $g$ در نقاط لازم مشتق‌پذیر باشند،
\[
\frac{d}{dt}g(u(t))=g'(u(t))u'(t).
\]
در کاربرد ما، تابع درونی آفین است. مثلاً با ثابت‌بودن $x,b$،
\[
\frac{\partial}{\partial w_i}g(w^{\mathsf T}x+b)
=x_i g'(w^{\mathsf T}x+b).
\]
برای دیدن مستقیم این حالت، خارج‌قسمتِ تفاضلی را در $w_i$ بنویسید و تغییرِ آرگومان را $rx_i$ بگیرید. اگر $x_i\neq0$ باشد، تعریف مشتقِ $g$ عامل $x_i$ را بیرون می‌آورد؛ اگر $x_i=0$ باشد، هر دو طرف صفرند. تکرار همین محاسبه فرمول مشتق‌های بالاتر را می‌دهد.

\begin{intuition}
مشتق‌گیری نسبت به وزن $w_i$، عامل $x_i$ را بیرون می‌آورد؛ مشتق‌گیری نسبت به خودِ ورودی $x_i$، عامل $w_i$ را بیرون می‌آورد. این دو را نباید جابه‌جا کرد.
\end{intuition}

\clearpage
\ribbon{پیوست پ | ادامه}
\heading{قضیهٔ مقدار میانگین و نتیجهٔ مشتق صفر}
اگر $v$ روی $[a,b]$ پیوسته و در $(a,b)$ مشتق‌پذیر باشد، نقطه‌ای $c\in(a,b)$ وجود دارد که
\[
v'(c)=\frac{v(b)-v(a)}{b-a}.
\]
\textbf{اثبات کوتاه.} خطِ وصل‌کنندهٔ دو سر نمودار را از $v$ کم کنید؛ تابع حاصل در دو انتها مقدار برابر دارد. اگر ثابت نباشد، یکی از بیشینه یا کمینه‌های آن در داخل بازه رخ می‌دهد. مشتق در یک اکسترمم داخلی صفر است، چون خارج‌قسمت‌های تفاضلی از دو سمت علامت‌های مخالف دارند. پس مشتقِ تابع حاصل در یک نقطه صفر می‌شود و فرمول بالا به دست می‌آید.

\subhead{نتیجهٔ اول: مشتق صفر یعنی تابع ثابت}
اگر $v'=0$ در تمام یک بازه باشد، قضیه را بین هر دو نقطهٔ آن بازه اعمال می‌کنیم. اختلاف مقدارهای تابع صفر می‌شود؛ پس $v$ ثابت است.

\subhead{نتیجهٔ دوم: مشتق مرتبهٔ $k$ صفر یعنی درجهٔ کمتر از $k$}
اگر $v\in C^k(\R)$ و $v^{(k)}\equiv0$ باشد، $v$ یک چندجمله‌ای با درجهٔ حداکثر $k-1$ است. نماد $C^k$ یعنی مشتق‌های تا مرتبهٔ $k$ وجود دارند و پیوسته‌اند.

\textbf{اثبات.} برای $k=1$ همان نتیجهٔ اول را داریم. برای $k>1$، مشتق $v^{(k-1)}$ صفر است؛ پس $v^{(k-1)}=c$ برای یک ثابت $c$. از $v$ تابعِ $ct^{k-1}/(k-1)!$ را کم کنید. مشتق مرتبهٔ $k-1$ تابع باقی‌مانده صفر است. با استقرا، باقی‌مانده درجهٔ حداکثر $k-2$ دارد؛ بنابراین درجهٔ $v$ حداکثر $k-1$ است.

\begin{intuition}
برای نمونه، اگر $v^{(3)}$ در همه‌جا صفر باشد، $v$ شکلی مانند $At^2+Bt+C$ دارد. صفر بودن مشتق فقط در یک نقطه چنین نتیجه‌ای نمی‌دهد؛ در گام ۳ اثبات اصلی، صفر بودن در \textbf{تمام خط} فرض شده بود.
\end{intuition}

\subhead{کاربرد مستقیم در میانگین‌گیری}
از صورت اولِ قضیهٔ اساسی حساب،
\[
\frac{d}{dt}\left(\frac1{2h}\int_{t-h}^{t+h}v(s)\,ds\right)
=\frac{v(t+h)-v(t-h)}{2h}.
\]
سمت راست فقط به مقدارهای $v$ وابسته است، نه به مشتق آن. بنابراین برای ساختن یک مشتق پیوسته از میانگین، پیوستگیِ خود $v$ کافی است. همین نکته اجازه می‌دهد اثبات برای فعال‌سازی‌هایی مثل ReLU هم کار کند.
\backtomain{step:three}

\appstart{ت}{درون‌یابی لاگرانژ و حفظِ سقف درجه}{app:lagrange}
\textbf{محل استفاده:} گام ۳ و اثباتِ کافی نبودن فعال‌سازی چندجمله‌ای.

\subhead{قضیهٔ درون‌یابی لاگرانژ}
برای $m+1$ نقطهٔ متمایز $t_0,\ldots,t_m$ و مقدارهای دلخواه $y_0,\ldots,y_m$، یک چندجمله‌ای یکتا با درجهٔ حداکثر $m$ وجود دارد که $p(t_i)=y_i$. فرمول آن
\[
p(t)=\sum_{i=0}^{m}y_i L_i(t),
\qquad L_i(t)=\prod_{j\neq i}\frac{t-t_j}{t_i-t_j}
\]
است. اگر $m=0$ باشد، حاصل‌ضرب تهی را یک می‌گیریم.

\textbf{اثبات.} هر $L_i$ در نقطهٔ $t_i$ یک و در بقیهٔ نقاط انتخاب‌شده صفر است؛ پس فرمول، مقدارهای خواسته‌شده را می‌گیرد. اختلافِ دو جواب احتمالی، درجهٔ حداکثر $m$ و دست‌کم $m+1$ ریشهٔ متمایز دارد؛ در نتیجه صفر است. دلیلِ محدودیت تعداد ریشه‌ها نیز ساده است: هر ریشهٔ متمایز یک عامل خطی جدا ایجاد می‌کند و درجه نمی‌تواند از تعداد این عامل‌ها کمتر باشد.

\begin{intuition}
یک خط با مقدارش در دو نقطه تعیین می‌شود؛ یک چندجمله‌ایِ درجهٔ حداکثر دو با مقدارش در سه نقطه. این قضیه همان ایده را برای درجهٔ دلخواه بیان می‌کند.
\end{intuition}

\subhead{لمِ موردنیاز: حدِ درجه‌های محدود}
اگر همهٔ $p_n$ها درجهٔ حداکثر $m$ داشته باشند و روی یک بازهٔ با طول مثبت به $u$ نقطه‌ای همگرا شوند، $u$ نیز روی آن بازه برابر یک چندجمله‌ایِ درجهٔ حداکثر $m$ است. بازه می‌تواند تمام $\R$ باشد.

\textbf{اثبات.} تعداد $m+1$ نقطهٔ متمایز را در بازه ثابت انتخاب کنید. برای هر $n$،
\[
p_n(t)=\sum_{i=0}^{m}p_n(t_i)L_i(t).
\]
با حدگیری در یک جمع متناهی،
\[
u(t)=\sum_{i=0}^{m}u(t_i)L_i(t).
\]
سمت راست، مستقل از $n$ و چندجمله‌ایِ درجهٔ حداکثر $m$ است. حتی همگرایی نقطه‌ای برای این نتیجه کافی بود؛ در اثبات اصلی، همگرایی یکنواخت نیز داریم.

\begin{intuition}
این لم نمی‌گوید حدِ هر دنباله‌ای از چندجمله‌ای‌ها چندجمله‌ای است. شرط مهم، \textbf{سقف ثابت برای درجهٔ همهٔ اعضا} است. در قضیهٔ وایرشتراس، درجه‌ها می‌توانند افزایش پیدا کنند؛ بنابراین تناقضی وجود ندارد.
\end{intuition}
\backtomain{step:three}

\appstart{ث}{لم هموارسازی و مجموع‌های ریمان}{app:smoothing}
\textbf{محل استفاده:} گام ۲؛ جایگزین‌کردن فعال‌سازیِ صرفاً پیوسته با یک ابزار مشتق‌پذیر.

\subhead{صورت لم}
برای $\sigma\in C(\R)$، عدد صحیح ثابت $k\geq1$ و $h>0$ بگذارید
\[
(A_hv)(t)=\frac1{2h}\int_{t-h}^{t+h}v(s)\,ds,
\qquad g_h=A_h^k\sigma.
\]
سه ویژگی داریم: $g_h\in C^k(\R)$؛ با $h\to0$ روی هر بازهٔ فشرده به $\sigma$ یکنواختاً نزدیک می‌شود؛ و برای هر $w,b$، تابع $g_h(w^{\mathsf T}x+b)$ در $\V$ قرار دارد.

\subhead{اثبات ویژگی اول: مشتق‌پذیری}
از پیوست پ،
\[
(A_hv)'(t)=\frac{v(t+h)-v(t-h)}{2h}.
\]
اگر $v\in C^r$ باشد، سمت راست $r$ مشتق پیوسته دارد؛ پس $A_hv\in C^{r+1}$. با $k$ بار تکرار، از پیوستگی $\sigma$ به $g_h\in C^k$ می‌رسیم.

\subhead{اثبات ویژگی دوم: نزدیک‌شدن به فعال‌سازی اصلی}
بازکردنِ $k$ میانگین‌گیری می‌دهد
\begin{equation}\label{eq:repeated-integral}
g_h(t)=\frac1{(2h)^k}\int_{[-h,h]^k}
\sigma(t+s_1+\cdots+s_k)\,ds_1\cdots ds_k.
\end{equation}
همهٔ وزن‌ها مثبت‌اند و میانگینِ تابع ثابتِ یک برابر یک است. همچنین $|s_1+\cdots+s_k|\leq kh$. بنابراین
\[
|g_h(t)-\sigma(t)|\leq
\sup_{|u|\leq kh}|\sigma(t+u)-\sigma(t)|.
\]
روی $|t|\leq R$ و برای $kh\leq1$، تمام آرگومان‌ها در بازهٔ ثابت $[-R-1,R+1]$ قرار می‌گیرند. پیوستگی یکنواخت $\sigma$ روی این بازه نشان می‌دهد سمت راست، به‌طور یکنواخت نسبت به $t$، به صفر می‌رود.

\subhead{نمونه: میانگین‌گیری از ReLU}
اگر $\sigma(t)=\max(0,t)$، یک بار میانگین‌گیری می‌دهد
\[
(A_h\sigma)(t)=
\begin{cases}
0,&t\leq-h,\\
(t+h)^2/(4h),&-h<t<h,\\
t,&t\geq h.
\end{cases}
\]
شکستگیِ صفر نرم می‌شود و مشتق در مرزها نیز پیوسته است. برای درجهٔ بالاتر، میانگین‌گیری را تکرار می‌کنیم؛ یک بار میانگین‌گیری، همهٔ مشتق‌ها را فراهم نمی‌کند.

\clearpage
\ribbon{پیوست ث | ادامه}
\heading{ویژگی سوم: چرا شبکهٔ اصلی حفظ می‌شود؟}
\subhead{لمِ مجموع‌های ریمانِ یکنواخت}
فرض کنید $D$ یک مستطیل بسته و کراندار و $\Phi(x,s)$ روی $K\times D$ پیوسته باشد. مجموع‌های ریمانِ انتگرالِ $\Phi$ نسبت به $s$، با ریزشدن افراز، \emph{یکنواخت نسبت به $x\in K$} همگرا می‌شوند.

\textbf{اثبات.} زیرانتگرال روی $K\times D$ یکنواختاً پیوسته است. افراز را آن‌قدر ریز کنید که برای هر $x$، اختلاف مقدارهای $\Phi(x,\cdot)$ در هر سلول کمتر از $\eta$ باشد. اگر $s_j$ نقطهٔ نمونهٔ سلول $D_j$ باشد،
\[
\left|\int_D\Phi(x,s)\,ds-
\sum_j\Phi(x,s_j)\operatorname{vol}(D_j)\right|
\leq\eta\operatorname{vol}(D).
\]
این کران به $x$ وابسته نیست. پس با کوچک‌شدن $\eta$، همگرایی یکنواخت حاصل می‌شود. همین کنترلِ نوسان، وجود انتگرال ریمانِ تابع پیوسته را هم توجیه می‌کند.

برای تابع پیوسته روی یک مستطیل، انتگرال تکراری با انتگرال روی مستطیل برابر است. کافی است افراز حاصل‌ضربی بگیریم: جمع‌های متناهی را می‌توان به هر ترتیب انجام داد، و همگرایی یکنواخت اجازه می‌دهد حد را در هر انتگرال بعدی عبور دهیم؛ زیرا
\[
\left|\int_a^b(u_n-u)\right|\leq(b-a)\sup_{[a,b]}|u_n-u|.
\]
در این کاربرد، به نظریهٔ اندازه یا صورت عمومی قضیهٔ فوبینی نیاز نداریم.

\subhead{کاربرد لم در شبکه}
در \eqref{eq:repeated-integral} به‌جای $t$ مقدار $w^{\mathsf T}x+b$ بگذارید. مجموع ریمانِ نرمال‌شده شکلی مانند
\[
\sum_{j=1}^{M}c_j\sigma(w^{\mathsf T}x+\widetilde b_j),
\qquad \widetilde b_j=b+s_{j,1}+\cdots+s_{j,k}
\]
دارد. این یک شبکهٔ متناهی با همان فعال‌سازی $\sigma$ است. لم نشان می‌دهد این شبکه‌ها روی $K$ یکنواختاً به $g_h(w^{\mathsf T}x+b)$ نزدیک می‌شوند. پس تابعِ موردنظر در $\V$ است.

\begin{intuition}
نمی‌گوییم «انتگرال یک شبکهٔ بی‌نهایت‌نورونی است» و کار را تمام کنیم. مجموع‌های ریمان دقیقاً نشان می‌دهند برای هر دقت، تعداد \textbf{متناهی} نورون و یک مجموعه پارامترِ مشترک برای همهٔ ورودی‌ها کافی است.
\end{intuition}

\subhead{یک مرز مهم}
هر انتخابِ $h$ لزوماً مفید نیست. برای نمونه، میانگین $\sin t$ روی یک دوره صفر می‌شود. گام ۳ فقط وجودِ یک $h$ و یک $b$ مناسب را برای هر درجه می‌خواهد؛ نه اینکه تمام هموارسازی‌ها غیرچندجمله‌ای بمانند.
\backtomain{step:two}

\appstart{ج}{چرا مشتق نسبت به وزن در مجموعهٔ قابل تقریب می‌ماند؟}{app:derivative}
\textbf{محل استفاده:} گام ۴؛ بخش فنیِ استخراج تک‌جمله‌ای‌های چندمتغیره.

\subhead{صورت لم}
فرض کنید $g\in C^k(\R)$، بایاس $b$ ثابت و $g(w^{\mathsf T}x+b)\in\V$ برای \emph{هر} وزن $w$ باشد. آنگاه هر مشتقِ متوالی تا مرتبهٔ $k$ نسبت به وزن‌ها، به‌عنوان تابعی از $x$، در $\V$ قرار می‌گیرد.

\subhead{اثباتِ مشتق اول}
بردار $e_i$ را برداری بگیرید که مؤلفهٔ $i$ام آن یک و بقیه صفر است. برای $r\neq0$، خطی‌بودن $\V$ می‌دهد
\[
Q_r(x)=\frac{g((w+re_i)^{\mathsf T}x+b)-g(w^{\mathsf T}x+b)}{r}\in\V.
\]
با استفاده از رابطهٔ \eqref{eq:ftc-difference}،
\[
Q_r(x)=x_i\int_0^1g'(w^{\mathsf T}x+b+\theta r x_i)\,d\theta.
\]
نشان می‌دهیم حد به‌صورت یکنواخت برابر $x_i g'(w^{\mathsf T}x+b)$ است. عدد $M\geq1$ را چنان انتخاب کنید که $|x_i|\leq M$ برای تمام مختصات و تمام $x\in K$. برای وزنِ ثابت $w$ و $|r|\leq1$، همهٔ آرگومان‌های $g'$ در یک بازهٔ فشردهٔ ثابت قرار دارند.

روی این بازه $g'$ یکنواختاً پیوسته است. برای هر $\eta>0$، اگر $|r|M$ به‌اندازهٔ کافی کوچک باشد، اختلاف مقدارِ $g'$ در آرگومانِ جابه‌جاشده و آرگومان اصلی کمتر از $\eta/M$ است. از فرمول انتگرالی نتیجه می‌شود خطای $Q_r$ در \emph{تمام} $x\in K$ کمتر از $\eta$ است. پس حد یکنواخت و، با قاعدهٔ دوم پیوست الف، عضو $\V$ است.

\subhead{تکرار استدلال تا مرتبهٔ $k$}
برای استقرا، فرض کنید پس از $j<k$ مشتق، برای هر $w$، تابع
\[
P(x)\,g^{(j)}(w^{\mathsf T}x+b)\in\V
\]
باشد؛ $P(x)$ حاصل‌ضرب همان $j$ مختصه‌ای است که از مشتق‌ها بیرون آمده‌اند. در آغاز $j=0$ و $P=1$ است. خارج‌قسمتِ تفاضلیِ همین خانواده نسبت به وزن $w_i$ در $\V$ است و به
\[
P(x)x_i g^{(j+1)}(w^{\mathsf T}x+b)
\]
میل می‌کند. برای یکنواختی، همان استدلال قبلی را به $g^{(j)}$ اعمال می‌کنیم. عامل $P(x)x_i$ روی $K$ حداکثر $M^{j+1}$ است و مشتقِ بعدی پیوسته است؛ بنابراین خطا را می‌توان برای تمام $x$ کوچک کرد. این استقرا تا مرتبهٔ $k$ معتبر است.

با $r_i$ بار مشتق نسبت به $w_i$ و $\sum_i r_i=k$، حاصل برابر $x_1^{r_1}\cdots x_d^{r_d}g^{(k)}(w^{\mathsf T}x+b)$ است. این همان فرمولِ استفاده‌شده در گام ۴ است.
\backtomain{step:four}

\appstart{چ}{وایرشتراس؛ چرا چندجمله‌ای‌ها کافی‌اند؟}{app:weierstrass}
\textbf{محل استفاده:} پیش‌نیازِ مستقلِ اثبات استون–وایرشتراس در پیوست ح.

\subhead{صورت قضیه}
برای هر تابع پیوستهٔ $f:[a,b]\to\R$ و هر $\eta>0$، چندجمله‌ای $p$ وجود دارد که $\sup_{[a,b]}|f-p|<\eta$. درجهٔ $p$ می‌تواند با دقت افزایش یابد. تابع $f$ لازم نیست مشتق داشته باشد یا دارای سری تیلور باشد.

\subhead{اثبات با چندجمله‌ای‌های برنشتاین}
با تغییر متغیر آفین، کافی است بازهٔ $[0,1]$ را بررسی کنیم. برای $n\geq1$ تعریف کنید
\[
B_nf(x)=\sum_{j=0}^{n}f(j/n)\,b_{n,j}(x),
\qquad b_{n,j}(x)=\binom nj x^j(1-x)^{n-j}.
\]
این تابع یک چندجمله‌ای است. از بسط دوجمله‌ای و رابطهٔ $j\binom nj=n\binom{n-1}{j-1}$ به دست می‌آید
\[
\sum_j b_{n,j}=1,\qquad
\sum_j j b_{n,j}=nx,\qquad
\sum_j j(j-1)b_{n,j}=n(n-1)x^2.
\]
بنابراین، با $j^2=j(j-1)+j$،
\begin{equation}\label{eq:bernstein}
\sum_j (j/n-x)^2b_{n,j}(x)=\frac{x(1-x)}n\leq\frac1{4n}.
\end{equation}
وزن‌ها نامنفی‌اند و مجموعشان یک است. رابطهٔ آخر می‌گوید وزنِ نقطه‌های دور از $x$ با بزرگ‌شدن $n$ کم می‌شود.

عدد $\eta>0$ را ثابت کنید. از پیوستگی یکنواخت، $\delta>0$ انتخاب کنید که فاصلهٔ کمتر از $\delta$، اختلافِ مقدار تابع کمتر از $\eta/2$ بدهد. بگذارید $M=\sup_{[0,1]}|f|$. خطا را به نمونه‌های نزدیک و دور تقسیم می‌کنیم:
\[
|B_nf(x)-f(x)|\leq\sum_j|f(j/n)-f(x)|\,b_{n,j}(x).
\]
سهم نمونه‌های با $|j/n-x|<\delta$ حداکثر $\eta/2$ است. برای بقیه، اختلافِ تابع حداکثر $2M$ و از \eqref{eq:bernstein}، مجموع وزن‌ها حداکثر $1/(4n\delta^2)$ است. پس
\[
\sup_{[0,1]}|B_nf-f|\leq\frac\eta2+\frac{M}{2n\delta^2}.
\]
برای $n$ کافی بزرگ، خطا کمتر از $\eta$ می‌شود. اگر $M=0$ باشد، خود تابع صفر است و نتیجه بدیهی است. تغییر متغیر، حکم را به هر بازهٔ غیرتک‌نقطه‌ای می‌برد؛ بازهٔ تک‌نقطه‌ای با یک ثابت حل می‌شود.

\begin{intuition}
این ساخت از مقدارهای تابع استفاده می‌کند، نه از مشتق‌های آن. بنابراین برای تابعی مثل $|x|$ هم معتبر است. هیچ بخش این برهان به قضیهٔ شبکهٔ عصبی وابسته نیست؛ در استدلال دور منطقی نداریم.
\end{intuition}
\backtomain{step:five}

\appstart{ح}{استون–وایرشتراس و تقریب چندمتغیره}{app:stone}
\textbf{محل استفاده:} گام ۵؛ عبور از چندجمله‌ای‌ها به تابع هدف روی $K\subset\R^d$.

\subhead{صورت حقیقیِ موردنیاز}
فرض کنید $K\subset\R^d$ فشردهٔ ناتهی و $\mathcal A$ خانواده‌ای از توابع پیوستهٔ حقیقی روی آن باشد. اگر این خانواده شامل ثابت‌ها باشد، تحت جمع و ضرب در عدد و ضربِ توابع بسته باشد، و نقاط را جدا کند، هر تابع پیوسته روی $K$ را یکنواختاً تقریب می‌زند.

«جداکردن نقاط» یعنی برای هر $x\neq y$، تابعی $a\in\mathcal A$ وجود داشته باشد که $a(x)\neq a(y)$. خانوادهٔ بسته تحت عملیات بالا را \textbf{جبر} می‌نامند.

\subhead{چرا در اثبات ما قابل استفاده است؟}
محدودیتِ چندجمله‌ای‌های چندمتغیره به $K$، ثابت‌ها را دارد و تحت جمع و ضرب بسته است. اگر دو نقطه متفاوت باشند، دست‌کم یکی از مختصاتشان فرق دارد؛ تابعِ مختصاتی $p(x)=x_i$ آن‌ها را جدا می‌کند. پس همهٔ شرط‌های قضیه برقرارند.

\begin{intuition}
در اثبات اصلی فقط همین صورتِ قضیه و بررسیِ شرط‌ها لازم است. برهان زیر برای آشنایی با خودِ ابزار آمده است و از نتیجهٔ تقریب شبکه‌ها استفاده نمی‌کند.
\end{intuition}

\subhead{اثبات ــ بخش اول: ساختن کمینه و بیشینه}
بگذارید $\mathcal B$ مجموعهٔ توابعی باشد که با اعضای $\mathcal A$ یکنواختاً قابل تقریب‌اند. طبق پیوست الف، این مجموعه خطی و نسبت به حد یکنواخت بسته است. تحت ضرب نیز بسته است: اگر $u_n\to u$ و $v_n\to v$ یکنواخت و $u_n,v_n\in\mathcal A$ باشند، $\|u_n\|_\infty$ کراندار است و
\[
\|u_nv_n-uv\|_\infty
\leq\|u_n\|_\infty\|v_n-v\|_\infty+
\|v\|_\infty\|u_n-u\|_\infty\longrightarrow0.
\]
پس $uv\in\mathcal B$.

برای $u\in\mathcal B$، عدد $M>\|u\|_\infty$ بگیرید. طبق وایرشتراسِ پیوست چ، چندجمله‌ای‌هایی $p_n(t)$ روی $[-M,M]$ به $|t|$ یکنواختاً نزدیک می‌شوند. چون $\mathcal B$ جبر است، $p_n(u)\in\mathcal B$. حدِ آن‌ها $|u|$ است؛ پس $|u|\in\mathcal B$.

از دو هویت زیر نتیجه می‌شود کمینه و بیشینهٔ متناهیِ اعضای $\mathcal B$ نیز در آن قرار دارند:
\[
\max(u,v)=\frac{u+v+|u-v|}{2},
\qquad \min(u,v)=\frac{u+v-|u-v|}{2}.
\]
این ابزار اجازه می‌دهد تقریب‌های محلی را کنار هم بگذاریم.

\clearpage
\ribbon{پیوست ح | ادامهٔ اثبات استون–وایرشتراس}
\heading{بخش دوم: از دو نقطه به کل دامنه}
تابع پیوستهٔ هدف $f$ روی $K$ و دقت $\eta>0$ را ثابت و $\delta=\eta/3$ قرار دهید.

برای هر دو نقطهٔ متفاوت $x,y$، تابع $a\in\mathcal A$ را با $a(x)\neq a(y)$ انتخاب کنید. تابع
\[
h_{x,y}(z)=f(x)+\frac{f(y)-f(x)}{a(y)-a(x)}\bigl(a(z)-a(x)\bigr)
\]
عضو $\mathcal A$ است و در هر دو نقطه مقدارِ $f$ را می‌گیرد. برای $x=y$، تابع ثابتِ $f(x)$ را به کار می‌بریم.

\subhead{نخست $x$ را ثابت نگه می‌داریم}
با تغییر $y$، مجموعه‌های باز
\[
U_y=\{z\in K:h_{x,y}(z)>f(z)-\delta\}
\]
کل $K$ را می‌پوشانند، چون $y\in U_y$. زیرپوشش متناهی با نقاط $y_1,\ldots,y_m$ انتخاب و تعریف کنید
\[
g_x(z)=\max_{1\leq i\leq m}h_{x,y_i}(z).
\]
این تابع در $\mathcal B$ است، همه‌جا از $f-\delta$ بزرگ‌تر است، و در خودِ $x$ دقیقاً با $f(x)$ برابر می‌شود.

\subhead{سپس $x$ را تغییر می‌دهیم}
مجموعهٔ باز $V_x=\{z:g_x(z)<f(z)+\delta\}$ شامل $x$ است. این مجموعه‌ها پوشش $K$ هستند. زیرپوشش متناهی با نقاط $x_1,\ldots,x_r$ انتخاب و قرار دهید
\[
g(z)=\min_{1\leq j\leq r}g_{x_j}(z).
\]
باز هم $g\in\mathcal B$ است. همهٔ $g_{x_j}$ها از $f-\delta$ بالاترند و در هر نقطه دست‌کم یکی از آن‌ها از $f+\delta$ پایین‌تر است. بنابراین
\[
f-\delta<g<f+\delta\quad\hbox{روی }K,
\qquad \|f-g\|_\infty\leq\delta.
\]
چون $g\in\mathcal B$، عضوی $a\in\mathcal A$ با $\|g-a\|_\infty<\delta$ وجود دارد. در نتیجه
\[
\|f-a\|_\infty<2\delta<\eta.
\]
قضیهٔ استون–وایرشتراس ثابت شد.\hfill$\square$

\begin{intuition}
دو بار از فشردگی استفاده کردیم: یک بار برای انتخابِ تعداد متناهی تقریب که از پایین کافی باشند؛ بار دوم برای کنترل از بالا. گرفتن بیشینه و کمینه، این تقریب‌ها را به یک تابعِ نزدیک به هدف تبدیل کرد.
\end{intuition}
\backtomain{step:five}

\clearpage
\phantomsection\label{sec:references}
\pdfbookmark[0]{منابع و حدود نتیجه}{bm-refs}
\ribbon{منابع و حدود نتیجه}
\heading{این نتیجه چه چیزی را نمی‌گوید؟}
قضیه، وجودِ وزن‌های مناسب را نشان می‌دهد. از آن، موفقیتِ گرادیان‌نزولی، تعداد کم نورون‌ها، تعمیم از دادهٔ محدود یا پایداری عددیِ ساختِ ارائه‌شده نتیجه نمی‌شود. ضرایبِ تفاضل‌ها ممکن است بزرگ باشند. همچنین بحثِ ضرورت برای \textbf{یک لایهٔ پنهان} و تقریب روی \textbf{تمام فشرده‌ها} است؛ دربارهٔ عمقِ قابل افزایش حکم تازه‌ای نداده‌ایم.

\heading{انتساب علمی و تفاوت با متن مقاله}
مرجع اصلی، قضیهٔ ۱ و بخش ۶ مقالهٔ Leshno و همکاران \cite{leshno1993} است. استخراج توان‌های ورودی با مشتق نسبت به وزن در گام سوم، صفحهٔ ۸۶۴ آن مقاله آمده است. مقاله کلاس عمومی‌تری از فعال‌سازی‌ها را بررسی می‌کند؛ این سند فقط حالت پیوسته را اثبات می‌کند.

در مقالهٔ مروری Pinkus \cite{pinkus1999}، صورت پیوسته در قضیهٔ \lr{3.1} و مسیر مشتق‌گیری و هموارسازی در گزاره‌های \lr{3.4} و \lr{3.7}، به‌ویژه صفحات ۱۵۶ تا ۱۵۸، توضیح داده شده است. این مقاله منبع نخستِ نتیجه نیست.

در بازنویسی حاضر، هموارسازی با میانگین‌گیری تکراری روی بازه انجام شده؛ حفظ درجه با درون‌یابی لاگرانژ نشان داده شده؛ و گذار چندمتغیره مستقیماً با مشتق نسبت به وزن‌ها نوشته شده است. این تغییرها ابزار توضیح‌اند، نه ادعای نتیجه‌ای جدید یا نقل لفظ‌به‌لفظ منابع.

\heading{مرز خودبسندگی}
لم‌های اختصاصیِ اثبات، همچنین برهان وایرشتراس و استون–وایرشتراس، در همین پیوست‌ها آمده‌اند. صورت هاینه–بورل به‌عنوان نتیجهٔ کلاسیک پذیرفته شده است. تعاریف حد، مشتق، انتگرال ریمان، کامل‌بودن اعداد حقیقی و جبر مقدماتی پیش‌نیاز مطالعه‌اند.

\renewcommand{\refname}{مراجع}
\begin{thebibliography}{9}
\begin{latin}
\bibitem{leshno1993}
M. Leshno, V. Ya. Lin, A. Pinkus, and S. Schocken.
\textit{Multilayer feedforward networks with a nonpolynomial activation function can approximate any function}.
Neural Networks, 6(6):861--867, 1993.
\href{https://doi.org/10.1016/S0893-6080(05)80131-5}{doi:10.1016/S0893-6080(05)80131-5}.
\newline
\href{https://pinkus.net.technion.ac.il/files/2021/02/neural.pdf}{Author-hosted full-text PDF}.

\bibitem{pinkus1999}
A. Pinkus.
\textit{Approximation theory of the MLP model in neural networks}.
Acta Numerica, 8:143--195, 1999.
\href{https://doi.org/10.1017/S0962492900002919}{doi:10.1017/S0962492900002919}.
\newline
\href{https://pinkus.net.technion.ac.il/files/2021/02/acta.pdf}{Author-hosted full-text PDF}.
\end{latin}
\end{thebibliography}
